\begin{abstract}
\noindent La presente relazione espone la \textbf{progettazione}, la \textbf{programmazione} e la \textbf{validazione} di un'unità di calcolo \textit{embedded} custom basata sul SoC \textit{Raspberry Pi RP2040}, ottimizzata per il \textit{data-logging} georeferenziato in contesti IoT. L'architettura proposta sfrutta il determinismo delle unità di I/O programmabili (\textit{Programmable I/O} - PIO) per l'offloading hardware di task critici di timing, preservando la reattività dei core \textit{ARM Cortex-M0+} durante la gestione delle periferiche di comunicazione. Particolare attenzione è rivolta alla robustezza del sistema, garantita dall'implementazione del \textit{filesystem fail-safe LittleFS} su memoria flash QSPI esterna e dall'adozione di rigorosi standard di manifattura SMT (\textit{Surface Mount Technology}). I risultati sperimentali dimostrano l'efficacia dell'integrazione tra sensori ambientali e moduli GNSS, confermando la resilienza dell'architettura in scenari di acquisizione dati ad alta frequenza e bassa latenza.
\end{abstract}

\vspace{0.5em}
\noindent \textbf{Keywords:} RP2040, SMT Manufacturing, GNSS Integration, PIO State Machines, IoT Resilience, Embedded Systems.
